% Session 5 -- The Report Number and the Memo: Procedures TPS-P-04, TPS-P-05.
% Law 1 sources: spec/tps_ritual.json R3 (1 to 6 digits, keyed twice; the
% two must agree or both are asked again; TPS-003) and R4 (the memo number
% must equal MEMO=; otherwise asked again; TPS-004); os/milton/tpsrep_
% program.asm .r3/.r4 and getfield (digits and the hyphen are accepted, at
% most 6 / 8 characters, further keys ignored, ENTER ends the entry, no
% echo except the accepted characters; an empty entry never agrees),
% rnoout (the report number zero-filled to six digits on the cover sheet and
% the page heads).  Screen tpsritual: the journey route.
% GATES: test-tps-report test-flair-tps-journey test-pubs-screens test-pubs-tps
\Gates{test-tps-report, test-flair-tps-journey, test-pubs-screens, test-pubs-tps}
\chapter{The Report Number and the Memo}

\SideHead{Introduction}When the form is aligned, the program asks for the
TPS report number, and then asks for it again.  When the two entries agree,
it asks for the number of the memo about the new cover sheets.  In this
session you key the report number twice, acknowledge the memo, and record
both in their logs.

\SideHead{Objectives}Upon completion of this session you will be able to:
\begin{enumerate}
\item Key the report number twice and recognize the refusal of numbers that
  do not agree.
\item Record the report number in the report number control log, \form{2}.
\item Key the memo number and recognize its refusal.
\item Record the memo acknowledgement in \form{3}.
\end{enumerate}

\section{The Report Number}
\index{report number}
The report number is assigned before the run and entered in \form{2},
column 1 (Checklist TPS-C-1).  It is one to six digits.  The program asks
for it with:
\TpsMsg{R3.prompt}
The first line is displayed first; when the first entry has been ended with
\key{Enter}, the second line is displayed.  Only the characters accepted are
displayed as they are typed.  After six characters, further keys are
ignored.  \key{Backspace} is ignored.

If the two entries do not agree, or are empty, the program displays:
\TpsMsg{TPS-003}
and asks for both entries again.  When they agree, the program continues
with the memo.  On the cover sheet and at the head of each page the report
number is printed with six digits, zeros first: report number 1014 is
printed \texttt{001014}.

\begin{Procedure}{TPS-P-04}{To Key the Report Number Twice}
\ProcField{Purpose}{To enter the assigned report number so that the program
  accepts it, and to record it.}
\ProcField{Frequency}{Each run, directly after \proc{03}.}
\ProcField{Performed by}{The operator.}
\ProcField{Prerequisites}{\TpsLine{R2.accept}{1} is displayed.  The report
  number is entered in \form{2}, column 1.}
\ProcField{Materials}{\form{2}.}
\begin{ProcSteps}
\item At \TpsLine{R3.prompt}{1} type the report number from \form{2},
  column 1, and press \key{Enter}.
\item At \TpsLine{R3.prompt}{2} type the report number again, from
  \form{2}, column 1, not from the screen.  Press \key{Enter}.
\item If \TpsLine{TPS-003}{1} is displayed, begin again at step 1.
  Make a tally mark in \form{2}, column 3.
\item When the memo prompt is displayed, write the report number as you
  keyed it the second time in \form{2}, column 2, and initial column 4.
\end{ProcSteps}
\ProcField{Verification}{The memo prompt follows the second entry.  Columns
  1 and 2 of \form{2} agree.}
\ProcField{Records retained}{\form{2}.}
\ProcNote{If you key a wrong digit, complete the entry and press \key{Enter}.
  The entries will not agree, and both are asked for again.}
\end{Procedure}

\OpsShot[fromtext={ENTER TPS REPORT NUMBER}, totext={RE-ENTER TPS REPORT NUMBER@2},
  marks={1=RE-ENTER TPS REPORT NUMBER: 1041;2=TPS-003;3=RE-ENTER TPS REPORT NUMBER@2},
  label=fig:r3, caption={Report number 1014 keyed, then 1041: refused.  Keyed
  twice again: accepted.}]{tpsritual}

\begin{CalloutLegend}
\LegendRow{1}{Second entry}{The second entry differs from the first.}
\LegendRow{2}{Refusal}{Both entries are asked for again.}
\LegendRow{3}{Agreement}{The entries agree.  The program continues with the
  memo.}
\end{CalloutLegend}

\section{The Memo About the New Cover Sheets}
\index{memo acknowledgement}
The program asks for the number of the memo about the new cover sheets:
\TpsMsg{R4.prompt}
The number keyed must be the number in the \texttt{MEMO=} line of the
parameter record, character for character.  Digits and the hyphen are
accepted, up to eight characters.  If the number is not the memo number, the
program displays:
\TpsMsg{TPS-004}
and asks again.  It continues to ask until the memo number is keyed.  When it
is, the program continues with the cover sheet (Session~6).

\begin{Procedure}{TPS-P-05}{To Acknowledge the Memo}
\ProcField{Purpose}{To confirm, by its number, that the memo about the new
  cover sheets has been read, and to record the acknowledgement.}
\ProcField{Frequency}{Each run, directly after \proc{04}.}
\ProcField{Performed by}{The operator.}
\ProcField{Prerequisites}{The memo prompt is displayed.  You have read the
  memo about the new cover sheets.}
\ProcField{Materials}{Your copy of the memo; \form{3}.}
\begin{ProcSteps}
\item Type the memo number exactly as it is printed on the memo, with its
  hyphen, and press \key{Enter}.
\item If \TpsLine{TPS-004}{1} is displayed, check the number on the memo
  and type it again.  If it is refused a second time, see your supervisor:
  the memo you hold is not the memo on file.  Leave the program waiting at
  the prompt.
\item When the cover sheet prompt is displayed, enter on \form{3} the memo
  number, the date and your initials.
\end{ProcSteps}
\ProcField{Verification}{The cover sheet prompt follows the memo number.}
\ProcField{Records retained}{\form{3}.}
\end{Procedure}

\OpsShot[fromtext={KEY THE NUMBER OF THE MEMO}, totext={KEY THE NUMBER OF THE MEMO@2},
  marks={1=TPS-004;2=KEY THE NUMBER OF THE MEMO@2}, label=fig:r4,
  caption={Memo number 88-113 refused; 88-114 accepted.}]{tpsritual}

\begin{CalloutLegend}
\LegendRow{1}{Refusal}{88-113 is not the number of the memo on file.}
\LegendRow{2}{Acceptance}{The memo number.  The program continues with the
  cover sheet.}
\end{CalloutLegend}

\section{Lab Project 4. The Report Number and the Memo}
Continue from Lab Project 3: the program displays the report number prompt.
\begin{enumerate}
\item Enter report number 1014 in \form{2}, column 1.
\item Type \type{1014} and press \key{Enter}; at the second prompt type
  \type{1041} and press \key{Enter}.  Compare your screen with
  Figure~\ref{fig:r3}.
\item Carry out \proc{04}.
\item At the memo prompt type \type{88-113} and press \key{Enter}.  Compare
  your screen with Figure~\ref{fig:r4}.
\item Carry out \proc{05} with the memo number \type{88-114}.
\item Continue with Lab Project 5 (Session~6) without leaving the program.
\end{enumerate}

\section{Study Questions}
\begin{enumerate}
\item How many times is the report number keyed?
\item From where is the second entry of the report number taken?
\item How is report number 52 printed on the cover sheet?
\item Where does the program find the memo number it compares with yours?
\item The memo number is refused twice.  What do you do?
\end{enumerate}
